% !TeX root = ../main.tex
\section{Typing rules and LETREC}
\label{sec:letrec}

\subsection{Typing rules}

Typing rules define which transformations are allowed on the values of a
program, through the functions that apply to each type. A rule has
premises above a line and a conclusion below it. The basic statement is a
\textbf{typing judgment}
\[
\Gamma \vdash \mathit{exp} : T
\]
read ``in the typing context $\Gamma$, the expression $\mathit{exp}$ has type
$T$''. The context $\Gamma$ is a map from variables to types: the types of the
variables in scope. The lecture's first rule is functional abstraction:
\begin{mathpar}
\inferrule*[right=(Abs)]{\Gamma, x : T_1 \vdash \mathit{exp} : T_2}{\Gamma \vdash \lambda x : T_1 .\, \mathit{exp} : T_1 \to T_2}
\end{mathpar}
If, in the context $\Gamma$ extended with $x$ of type $T_1$, the body
$\mathit{exp}$ has type $T_2$, then the function that takes $x$ and returns
$\mathit{exp}$ has the function type $T_1 \to T_2$.

\subsection{The language LETREC}

The lecture uses a small language, LETREC, from Friedman and
Wand~\cite[Ch.~4]{friedman2008essentials}:
\begin{itemize}
  \item it has numbers and boolean values;
  \item variables are introduced with \texttt{let};
  \item the only arithmetic expression is the difference \texttt{-(e1, e2)};
  \item there are \texttt{if} expressions and the test \texttt{zero?};
  \item new functions can be defined (\texttt{proc}, and \texttt{letrec} for
    recursive ones) and applied (\texttt{(f e)}).
\end{itemize}
The grammar of the lecture is
\begin{lstlisting}
Expression ::= letrec Identifier (Identifier) = Expression in Expression
             | (Expression Expression)
             | proc (Identifier) Expression
             | if Bool-exp then Expression else Expression
             | num | bool
             | - (Expression, Expression)
Bool-exp   ::= Expression Operator Expression | Identifier (Expression)
Operator   ::= < | > | <> | || | && | ==
\end{lstlisting}
and a program that doubles a number by recursion:
\begin{lstlisting}
letrec double(x)
  = if zero?(x) then 0 else -((double -(x,1)), -2)
in (double 6)
\end{lstlisting}
Subtracting $-2$ adds 2, so \texttt{(double 6)} is 12. The typed version of the
language adds types to the parameters and to the result of \texttt{letrec}:
\begin{lstlisting}
Expression ::= letrec Type Identifier (Identifier : Type) = Expression in Expression
             | proc (Identifier : Type) Expression
             | ...
Type       ::= int | boolean
\end{lstlisting}
\begin{lstlisting}
letrec
  int double (x : int) = if zero?(x) then 0 else -((double -(x,1)), -2)
in double
\end{lstlisting}
Typed declarations can come before the variable (\texttt{T var}, as in C) or
after it (\texttt{var : T}, as here). A function can be typed by its
parameters and result, $T_{\mathit{res}}\ \mathit{foo}(T_1\ p_1, \dots, T_n\ p_n)$
or $\mathit{foo}(p_1 : T_1, \dots, p_n : T_n) : T_{\mathit{res}}$, or with its full type
$\mathit{foo} : T_1 \to T_2 \to \dots \to T_{\mathit{res}}$; the first two forms mean the
third.

In Rascal, the syntax with optional types (used also for inference in
Section~\ref{sec:inference}) and the abstract syntax are:
\coderange{Rascal}{rascal/src/Letrec.rsc}{syntax}
\coderange{Rascal}{rascal/src/Letrec.rsc}{ast}
In the concrete syntax a function type is written \texttt{(int -> int)}. In the
abstract syntax it is \texttt{procType(intType(), intType())}.

\subsection{Typing rules of LETREC}

\begin{mathpar}
\inferrule*[right=(Const)]{ }{\Gamma \vdash n : \mathit{int}}
\and
\inferrule*[right=(Var)]{(x : T) \in \Gamma}{\Gamma \vdash x : T}
\and
\inferrule*[right=(Pred)]{\Gamma \vdash e : \mathit{int}}{\Gamma \vdash \mathtt{zero?}(e) : \mathit{bool}}
\and
\inferrule*[right=(Diff)]{\Gamma \vdash e_1 : \mathit{int} \\ \Gamma \vdash e_2 : \mathit{int}}{\Gamma \vdash \mathtt{-}(e_1, e_2) : \mathit{int}}
\and
\inferrule*[right=(Var-Intro)]{\Gamma \vdash e : T_1 \\ \Gamma, x : T_1 \vdash \mathit{body} : T_2}{\Gamma \vdash \mathtt{let}\ x = e\ \mathtt{in}\ \mathit{body} : T_2}
\and
\inferrule*[right=(If)]{\Gamma \vdash e_1 : \mathit{bool} \\ \Gamma \vdash e_2 : T \\ \Gamma \vdash e_3 : T}{\Gamma \vdash \mathtt{if}\ e_1\ \mathtt{then}\ e_2\ \mathtt{else}\ e_3 : T}
\and
\inferrule*[right=(Reduc)]{\Gamma \vdash \mathit{rator} : T_1 \to T_2 \\ \Gamma \vdash \mathit{rand} : T_1}{\Gamma \vdash (\mathit{rator}\ \mathit{rand}) : T_2}
\end{mathpar}
Each rule reads from the conclusion: to know the type of a difference, check
that both operands are integers; the result is an integer. The rule
\textsc{If} requires both branches to have the same type, which becomes the
type of the whole expression. \textsc{Reduc} (application) requires the
operator to be a function whose argument type is the type of the operand.

\paragraph{Functions.} The rule for \texttt{proc} is \textsc{Abs}:
\begin{mathpar}
\inferrule*[right=(Abs)]{\Gamma, x : T_1 \vdash \mathit{body} : T_2}{\Gamma \vdash \mathtt{proc}\ (x : T_1)\ \mathit{body} : T_1 \to T_2}
\end{mathpar}
It needs the type of $x$. In the typed language the programmer writes it; in
Section~\ref{sec:inference} the checker deduces it. For \texttt{letrec}, the
function is in scope in its own body, so the context for the body already
contains $f$:
\begin{mathpar}
\inferrule*[right=(Letrec)]{\Gamma, f : T_1 \to T_2, x : T_1 \vdash \mathit{fbody} : T_2 \\ \Gamma, f : T_1 \to T_2 \vdash \mathit{body} : T}{\Gamma \vdash \mathtt{letrec}\ T_2\ f\,(x : T_1) = \mathit{fbody}\ \mathtt{in}\ \mathit{body} : T}
\end{mathpar}

\paragraph{A derivation.} The rules combine into a tree, the typing
derivation. Figure~\ref{fig:derivation} shows the one for
\texttt{let x = 3 in -(x, 1)}.

\begin{figure}[h]
\centering
\begin{mathpar}
\inferrule*[right=(Var-Intro)]
  {\inferrule*[right=(Const)]{ }{\vdash 3 : \mathit{int}}
   \\
   \inferrule*[right=(Diff)]
     {\inferrule*[right=(Var)]{(x : \mathit{int}) \in \{x : \mathit{int}\}}{x : \mathit{int} \vdash x : \mathit{int}}
      \\
      \inferrule*[right=(Const)]{ }{x : \mathit{int} \vdash 1 : \mathit{int}}}
     {x : \mathit{int} \vdash \mathtt{-}(x, 1) : \mathit{int}}}
  {\vdash \mathtt{let}\ x = 3\ \mathtt{in}\ \mathtt{-}(x, 1) : \mathit{int}}
\end{mathpar}
\caption{Typing derivation of \texttt{let x = 3 in -(x, 1)}.}
\label{fig:derivation}
\end{figure}

\subsection{The function type-of}

The rules become a function that computes the type of an expression in an
environment. The lecture writes it as
\begin{lstlisting}
fun type-of(tenv, exp) {
  if exp==num && const-exp(num) then int
  if exp==var && var-exp(var) then tenv(var)
  if type-of(exp) == int then t(zero?-exp(exp), bool)
  if type-of(exp1) == int && type-of(exp2) == int then t(diff(exp1 exp2), int)
  if type-of(exp) == T1 && type-of(body) == T2 && tenv(var) then t(let-exp(var exp body), T2)
  if type-of(exp1)==bool && type-of(exp2)==T && type-of(exp3)==T then t(if-exp(exp1, exp2, exp3), T)
  if type-of(rator) == T1 -> T2 && type-of(rand) == T1 then t(call-exp(rator, rand), T2)
}
\end{lstlisting}
In Rascal there is one equation per rule, and each \texttt{checkEqual} is a
premise:
\coderange{Rascal}{rascal/src/Letrec.rsc}{typeof}
The tests type the program \texttt{double} as \texttt{int -> int}, its
application to 6 as \texttt{int}, and reject an \texttt{if} whose branches have
different types and the application of a number.

\paragraph{Type consistency.} The core of the checker is the comparison of
types. Given two types that an expression requires to be equal, the checker
continues if they are and reports an error otherwise:
\begin{lstlisting}
error check-equal(T1, T2, exp) {
  if T1 != T2 then
    error("Types did not match in expression %s: %s != %s", exp, name(T1), name(T2))
  else continue
}
\end{lstlisting}
The Rascal version is \texttt{checkEqual} above. The lecture also shows the
same checks written with TypePal, a Rascal library for type checkers, where
each construct \texttt{collect}s its constraints and a solver computes the
types.

\subsection{Type checking and type inference}

\paragraph{Type checking.} At compile time all the entities of the program
(values, variables, functions, objects) must have a type. The programmer
provides the types; the type checker deduces the types of the other
expressions and checks that they are consistent. In
\begin{lstlisting}[language=C]
int res;
int x = 3;
res = x / 2.0;
\end{lstlisting}
the types of \texttt{res} and \texttt{x} are given. The checker deduces that
\texttt{x / 2.0} is \texttt{float}, since dividing an \texttt{int} by a
\texttt{float} gives a \texttt{float}, and then has to decide whether a
\texttt{float} can be assigned to an \texttt{int} variable. Java rejects it; C
converts the value implicitly.

\paragraph{Type inference.} The programmer does not provide all the types, and
the checker infers them from their use. In Haskell,
\coderange{Haskell}{haskell/DoubleSmall.hs}{double}
has no type annotation. \texttt{x > 100} needs \texttt{x} to be something that
can be compared, and \texttt{x * 2} something that can be multiplied, so the
function goes from numbers to numbers. The lecture writes this as
$\mathit{Num} \to \mathit{Num}$; GHC's exact answer is
\texttt{(Ord a, Num a) => a -> a}: for any type \texttt{a} that can be compared
and has arithmetic. Section~\ref{sec:inference} shows how such types are
computed.

\subsection{Type environments}

Variables are typed in an environment, the same way the symbol table works
for names. The lecture's example is in OCaml:
\begin{lstlisting}
let someval = "hello";;
let otherval = "world" in let someval = 3 in
  Printf.printf "%d %s \n" someval otherval;;
\end{lstlisting}
The first line gives \texttt{someval : string}. The second creates a new
environment for \texttt{otherval} and, inside it, another one where a new
\texttt{someval} of type \texttt{int} hides the first. The program prints
\texttt{3 world}. Figure~\ref{fig:tenv} shows the three environments. The same
structure in Scheme:
\coderange{Scheme}{scheme/dynamic.scm}{env}

\begin{figure}[h]
\centering
\begin{tikzpicture}[env/.style={draw, rounded corners=2pt, align=left, font=\small, inner sep=4pt}]
  \node[env] (g1) at (0,0) {$\Gamma_1$\\ \texttt{someval : string}};
  \node[env] (g2) at (5.2,0) {$\Gamma_2$\\ \texttt{otherval : string}};
  \node[env] (g3) at (10.4,0) {$\Gamma_3$\\ \texttt{someval : int}};
  \draw[flecha] (g2) -- node[above, font=\footnotesize] {extends} (g1);
  \draw[flecha] (g3) -- node[above, font=\footnotesize] {extends} (g2);
\end{tikzpicture}
\caption{Nested type environments. In $\Gamma_3$, $\Gamma_3 \vdash \mathtt{someval} : \mathit{int}$ because the search stops at the innermost declaration.}
\label{fig:tenv}
\end{figure}

\subsection{Exercises}

\begin{exercise}
What would be the rule for the definition of functions in LETREC? How does the
checker know the type of the parameter?
\end{exercise}

\paragraph{Solution.} It is \textsc{Abs}: if the body has type $T_2$ when the
parameter has type $T_1$, the function has type $T_1 \to T_2$. The parameter's
type is written by the programmer in the typed language
(\texttt{proc (x : int) \dots}); without the annotation it has to be inferred
from how the body uses $x$.

\begin{exercise}
Give the typing derivation of
\texttt{letrec int double (x : int) = if zero?(x) then 0 else -((double -(x,1)), -2) in double}.
\end{exercise}

\paragraph{Solution.} With $\Gamma = \{\mathit{double} : \mathit{int} \to \mathit{int}\}$ and
$\Gamma' = \Gamma, x : \mathit{int}$. In $\Gamma'$: $\mathtt{zero?}(x) : \mathit{bool}$ by
\textsc{Pred} and \textsc{Var}; $0 : \mathit{int}$; $\mathtt{-}(x, 1) : \mathit{int}$ by
\textsc{Diff}; $(\mathit{double}\ \mathtt{-}(x,1)) : \mathit{int}$ by \textsc{Reduc}, since
$\mathit{double} : \mathit{int} \to \mathit{int}$; $\mathtt{-}(\dots, -2) : \mathit{int}$ by
\textsc{Diff}; so the \texttt{if} has type $\mathit{int}$ by \textsc{If}, which matches
the declared result. Then the body \texttt{double} has type $\mathit{int} \to \mathit{int}$
in $\Gamma$ by \textsc{Var}, and by \textsc{Letrec} the whole program has type
$\mathit{int} \to \mathit{int}$.

\begin{exercise}
Which rule fails for \texttt{if zero?(1) then 1 else zero?(0)}?
\end{exercise}

\paragraph{Solution.} \textsc{If}: the condition is a $\mathit{bool}$, but the branches
have types $\mathit{int}$ and $\mathit{bool}$, and the rule requires the same type
$T$ for both. The Rascal checker reports
\texttt{Types did not match \dots: bool != int}.

\begin{exercise}
In the OCaml example, what does \texttt{Printf.printf "\%d \%s" someval otherval}
print if the inner \texttt{let someval = 3} is removed? Does the program still
type check?
\end{exercise}

\paragraph{Solution.} It does not type check: \texttt{someval} is now the outer
one, a \texttt{string}, and \texttt{\%d} requires an \texttt{int}. The format
string makes the types of the arguments part of the check.
